% ==========================================
% N3 - GIAM TIMER BGP TREN L4
% ==========================================

\section{N3 -- Giảm timer BGP trên liên kết L4}

\subsection{Mục tiêu}

Trong F4, thời gian gián đoạn khi chỉ tắt
đầu A2 dài hơn nhiều so với khi tắt cả hai đầu.
Nhóm giảm timer BGP trên L4 để B1
phát hiện mất kết nối sớm hơn trong trường hợp
interface của nó vẫn ở trạng thái \texttt{up/up}.
Sau đó, nhóm lặp lại phép thử và so sánh với F4.

\subsection{Triển khai cấu hình}

Cấu hình ban đầu dùng Keepalive Interval
60 giây và Hold Time 180 giây.
Nhóm thay đổi hai giá trị này bằng lệnh:

\begin{lstlisting}
neighbor <IP-neighbor> timers 10 30
\end{lstlisting}

Trong đó, 10 giây là khoảng thời gian
gửi bản tin KEEPALIVE định kỳ;
30 giây là Hold Time đề xuất cho phiên BGP.
Nếu không nhận được bản tin KEEPALIVE hoặc UPDATE
trong khoảng Hold Time đã thương lượng,
router sẽ kết luận phiên không còn hoạt động.
Nhóm cấu hình trên cả hai đầu L4
và kiểm tra timer thực tế sau khi
phiên được thiết lập lại.

\textbf{Trên A2:}

\begin{lstlisting}
conf t
router bgp 64626
 neighbor 10.71.26.14 timers 10 30
end
\end{lstlisting}

\textbf{Trên B1:}

\begin{lstlisting}
conf t
router bgp 64726
 neighbor 10.71.26.13 timers 10 30
end
\end{lstlisting}

Nhóm dùng các lệnh sau để kiểm tra
trạng thái và timer của phiên:

\begin{lstlisting}
! A2
show ip bgp neighbors 10.71.26.14

! B1
show ip bgp neighbors 10.71.26.13
\end{lstlisting}

\partzeroimage
{images/advanced/N3/A2_bgp_timers.png}
{Timer của phiên eBGP tới B1 trên A2 sau khi thay đổi}
{fig:N3-a2-timers}

\partzeroimage
{images/advanced/N3/B1_bgp_timers.png}
{Timer của phiên eBGP tới A2 trên B1 sau khi thay đổi}
{fig:N3-b1-timers}

Cả hai output đều cho thấy phiên ở trạng thái
\texttt{Established}, \texttt{hold time is 30}
và \texttt{keepalive interval is 10 seconds}.
Đây là timer đang có hiệu lực trên phiên,
không chỉ là giá trị được ghi trong cấu hình.

\subsection{Thử lỗi và ghi nhận kết quả}

Nhóm giữ cách đo như F4:
ping liên tục từ A1 tới LAN-C,
dùng Loopback1 làm địa chỉ nguồn:

\begin{lstlisting}
ping 10.171.126.129 source Loopback1 repeat 100000 timeout 1
\end{lstlisting}

Khi mạng đã ổn định và ưu tiên đường qua B,
nhóm lần lượt thử chỉ tắt FastEthernet1/0
trên A2, rồi thử tắt cả cổng này và
FastEthernet0/0 trên B1.
Các lệnh \texttt{shutdown} giống F4.
Sau mỗi lần thử, nhóm bật lại các cổng
bằng \texttt{no shutdown}, chờ phiên eBGP
được thiết lập và mạng hội tụ trước lần tiếp theo.

\subsubsection{Chỉ tắt đầu A2}

\partzeroimage
{images/advanced/N3/one_end_failure_start.png}
{Mốc gây lỗi phía A2 sau khi giảm timer: 00:12:02.122}
{fig:N3-one-start}

\partzeroimage
{images/advanced/N3/one_end_ping_recovered.png}
{Mốc ping trở lại khi chỉ tắt đầu A2: 00:12:31.003}
{fig:N3-one-recovered}

Từ hai mốc trên đồng hồ bản ghi,
thời gian ping phục hồi là
\textbf{28,881 giây}.

\subsubsection{Tắt cả hai đầu L4}

\partzeroimage
{images/advanced/N3/both_ends_failure_start.png}
{Mốc bắt đầu thử lỗi ở cả hai đầu sau khi giảm timer: 00:14:03.433}
{fig:N3-both-start}

\partzeroimage
{images/advanced/N3/both_ends_ping_recovered.png}
{Mốc ping trở lại khi tắt cả hai đầu: 00:14:04.616}
{fig:N3-both-recovered}

Thời gian ping phục hồi trong lần thử này
là \textbf{1,183 giây}.

\subsection{So sánh với F4}

\begin{table}[H]
\centering
\caption{So sánh thời gian phục hồi trước và sau khi giảm timer BGP}
\label{tab:N3-comparison}
\small
\renewcommand{\arraystretch}{1.3}
\begin{tabularx}{\textwidth}{X c c}
\toprule
\textbf{Trường hợp} &
\makecell{\textbf{F4: timer 60/180}\\\textbf{Thời gian (s)}} &
\makecell{\textbf{N3: timer 10/30}\\\textbf{Thời gian (s)}} \\
\midrule
Chỉ tắt đầu A2 & 176,876 & 28,881 \\
Tắt cả hai đầu & 1,050 & 1,183 \\
\bottomrule
\end{tabularx}
\end{table}

Khi giảm Hold Time xuống 30 giây,
thời gian gián đoạn ở trường hợp chỉ tắt
một đầu giảm rõ rệt, từ 176,876 xuống
28,881 giây.
B1 có thể kết luận mất kết nối sớm hơn,
không phải chờ theo Hold Time 180 giây như trước.
Keepalive Interval 10 giây giúp bản tin
duy trì phiên được gửi thường xuyên hơn,
phù hợp với Hold Time ngắn hơn;
giá trị Hold Time mới là giới hạn chờ
để kết luận phiên mất hoạt động.

Thời gian phát hiện lỗi không nhất thiết
bằng đúng 30 giây tính từ lúc shutdown,
vì Hold Timer tính từ bản tin BGP hợp lệ
cuối cùng nhận được.
Phần thời gian còn lại của timer
phụ thuộc thời điểm gây lỗi.
Ngoài ra, số đo ping còn bao gồm
thời gian cập nhật tuyến và phục hồi
đường đi, đường về như đã phân tích
ở Mục \ref{sec:F4}.

Ở trường hợp tắt cả hai đầu,
hai router đều phát hiện interface của mình down
và có thể ngắt phiên ngay nhờ
Fast External Fallover,
không cần đợi Hold Time.
Vì vậy, giảm timer không tạo ra
khác biệt rõ rệt trong các lần đo này:
thời gian phục hồi vẫn xấp xỉ 1 giây.
Chênh lệch giữa 1,050 và 1,183 giây
chưa đủ để kết luận cấu hình mới chậm hơn,
vì đây là các lần đo riêng trên môi trường mô phỏng.

\newpage
